\section[Práctica con enteros y ecuaciones]{Ejercicios con números enteros y ecuaciones}

En los siguientes ejercicios utilizaremos conjuntamente las ideas desarrolladas hasta aquí. Recuerde que una sustracción puede escribirse como la suma del opuesto:
\[
a-b=a+(-b).
\]
También recuerde que, en una igualdad, cualquier transformación destinada a conservarla debe realizarse sobre ambos miembros.

\subsection*{Reescribir sustracciones}

Reescriba cada expresión utilizando únicamente adiciones. No es necesario evaluarlas todavía.
\begin{enumerate}
\item \(8-3\)

\item \(2-10\)

\item \(-5-7\)

\item \(14-(-5)\)

\item \(-8-(-3)\)

\item \(20-6-4\)

\item \(-15+8-(-3)\)

\item \(12-(-4)-9\)
\end{enumerate}
Por ejemplo,
\[
7-2+(-5)
\]
puede escribirse como
\[
7+(-2)+(-5).
\]
Una vez escrita de esta manera, resulta más sencillo identificar cada uno de los sumandos.

\subsection*{Evaluar expresiones con enteros}

Encuentre un único numeral que represente el valor de cada expresión.
\begin{enumerate}
\item \(12+(-7)\)

\item \((-9)+(-6)\)

\item \(8-15\)

\item \(14-(-5)\)

\item \(-8-(-3)\)

\item \(10-4+(-2)\)

\item \(-15+8-(-3)\)

\item \(45+(-28)-13+(-7)\)

\item \(-156+289-(-74)\)

\item \(-125+68-(-42)+(-89)-23\)
\end{enumerate}

Cuando resulte conveniente, reescriba primero todas las sustracciones como adiciones y reorganice luego los sumandos.

\subsection*{Completar igualdades}

Complete cada espacio con un número entero que haga verdadera la igualdad.

\begin{enumerate}
\item \(7+\square=12\)

\item \(15+\square=9\)

\item \(\square+8=3\)

\item \(10-\square=4\)

\item \(5-\square=12\)

\item \(-3+\square=8\)

\item \(12+(-5)=10+\square\)

\item \(20-8=\square+7\)

\item \(-4+9=12+\square\)

\item \(15-(-3)=20+\square\)
\end{enumerate}

Observe que aquí el objetivo no es simplemente realizar una operación: debe encontrar una cantidad que haga que ambos miembros representen el mismo número.

\subsection*{Resolver ecuaciones}

Encuentre el valor de \(x\) que hace verdadera cada ecuación. Escriba las transformaciones realizadas sobre ambos miembros.

\begin{enumerate}
\item

\[
x+4=11
\]

\item

\[
x+7=3
\]

\item

\[
x-5=12
\]

\item

\[
x-8=-2
\]

\item

\[
x+(-6)=10
\]

\item

\[
x-(-4)=9
\]

\item

\[
-7+x=5
\]

\item

\[
13+x=-4
\]

\item

\[
x+25=-18
\]

\item

\[
x-(-15)=-7
\]

\end{enumerate}

\subsection*{Comprobar soluciones}

Determine si el valor propuesto para \(x\) es solución de la ecuación. Para comprobarlo, sustituya \(x\) por el valor indicado y determine si obtiene una igualdad verdadera.

\begin{enumerate}
\item \(x+5=12\), con \(x=7\).

\item \(x-4=10\), con \(x=6\).

\item \(x+(-8)=-3\), con \(x=5\).

\item \(x-(-2)=9\), con \(x=7\).

\item \(-6+x=-10\), con \(x=-4\).

\item \(15+x=4\), con \(x=-11\).
\end{enumerate}
